\documentclass[tikz,border=10pt]{standalone}
\usepackage{ctex}
\usepackage{amsmath}
\usetikzlibrary{calc, arrows.meta, shapes.geometric, positioning, backgrounds, fit}

\begin{document}
% 定点 / 定值 / 定直线 三类问题对比示意图
\begin{tikzpicture}[
  font=\small,
  box/.style={draw, rounded corners=4pt, minimum width=3.4cm, minimum height=1.0cm,
              align=center, inner sep=6pt},
  arrow/.style={-Stealth, thick},
]

%--- 标题 ---
\node[font=\normalsize\bfseries] at (4.5, 6.8) {"3 定"高考压轴问题};

%─────────────────────────────────────
% 左栏：定点问题
%─────────────────────────────────────
\node[box, fill=blue!12, draw=blue!60] (fixpt-title) at (0.5, 5.5)
  {\textbf{定点问题}\\动线恒过某点};

% 示意图：动直线族过定点 P
\begin{scope}[shift={(0.5,3.0)}]
  % 椭圆
  \draw[blue!50, thick] (0,0) ellipse (1.1cm and 0.65cm);
  % 定点 P
  \coordinate (P) at (-0.3, 0.15);
  \node[fill=red, circle, inner sep=2pt] at (P) {};
  \node[red, above right, font=\scriptsize] at (P) {$P$（定点）};
  % 几条动直线（参数不同，但都过 P）
  \foreach \ang in {-60,-20,20,55}{
    \draw[blue!35, thin]
      ($(P)+({cos(\ang)*1.6},{sin(\ang)*1.6})$) --
      ($(P)-({cos(\ang)*1.6},{sin(\ang)*1.6})$);
  }
  \draw[blue!70, thick]
    ($(P)+({cos(20)*1.6},{sin(20)*1.6})$) --
    ($(P)-({cos(20)*1.6},{sin(20)*1.6})$);
\end{scope}

% 套路框
\node[draw=blue!40, fill=blue!5, rounded corners=3pt,
      font=\scriptsize, align=left, inner sep=5pt,
      text width=3.0cm] (fixpt-method) at (0.5, 1.3)
  {套路：\\
   ① 设参方程\\
   ② 收集参数项\\
   ③ 含参项 $=0$\\
   ④ 解出定点坐标};

%─────────────────────────────────────
% 中栏：定值问题
%─────────────────────────────────────
\node[box, fill=green!12, draw=green!60] (fixval-title) at (4.5, 5.5)
  {\textbf{定值问题}\\某量恒为常数};

\begin{scope}[shift={(4.5,3.0)}]
  % 椭圆
  \draw[green!60!black, thick] (0,0) ellipse (1.1cm and 0.65cm);
  % 弦 AB
  \coordinate (A) at (-0.9, 0.40);
  \coordinate (B) at (0.85, 0.40);
  \node[fill=black, circle, inner sep=1.5pt] at (A) {};
  \node[fill=black, circle, inner sep=1.5pt] at (B) {};
  \node[font=\scriptsize, above left] at (A) {$A$};
  \node[font=\scriptsize, above right] at (B) {$B$};
  \draw[green!60!black, thick] (A) -- (B);
  % 弦长标注
  \draw[<->, green!50!black, thin]
    ($(A)+(0,-0.25)$) -- ($(B)+(0,-0.25)$)
    node[midway, below, font=\scriptsize] {$|AB|=\text{定值}$};
\end{scope}

\node[draw=green!50, fill=green!5, rounded corners=3pt,
      font=\scriptsize, align=left, inner sep=5pt,
      text width=3.0cm] at (4.5, 1.3)
  {套路：\\
   ① 设弦端点参数\\
   ② 用韦达定理\\
   ③ 代入 $|AB|^2$\\
   ④ 化简，参数消去};

%─────────────────────────────────────
% 右栏：定直线问题
%─────────────────────────────────────
\node[box, fill=orange!12, draw=orange!70] (fixline-title) at (8.5, 5.5)
  {\textbf{定直线问题}\\中点轨迹为直线};

\begin{scope}[shift={(8.5,3.0)}]
  % 椭圆
  \draw[orange!70!black, thick] (0,0) ellipse (1.1cm and 0.65cm);
  % 若干弦的中点（连成直线）
  \foreach \yy/\xa/\xb in {
    0.45/-0.82/0.82,
    0.20/-1.05/1.05,
   -0.10/-1.08/1.08,
   -0.38/-0.90/0.90}{
    \coordinate (LL) at (\xa, \yy);
    \coordinate (RR) at (\xb, \yy);
    \coordinate (MM) at ({(\xa+\xb)/2}, \yy);
    \draw[orange!40, thin] (LL) -- (RR);
    \node[fill=red!80, circle, inner sep=1.3pt] at (MM) {};
  }
  % 中点连线（红色定直线）
  \draw[red, thick] (-0.2, -0.55) -- (0.2, 0.60)
    node[right, font=\scriptsize, red] {定直线};
\end{scope}

\node[draw=orange!50, fill=orange!5, rounded corners=3pt,
      font=\scriptsize, align=left, inner sep=5pt,
      text width=3.0cm] at (8.5, 1.3)
  {套路：\\
   ① 设弦中点 $(x_0,y_0)$\\
   ② 中点代入椭圆\\
      两点方程相减\\
   ③ 结合斜率关系\\
   ④ 得 $y_0=kx_0+b$};

%─────────────────────────────────────
% 共同箭头提示
%─────────────────────────────────────
\node[font=\scriptsize, gray, align=center] at (4.5, 0.4)
  {三类问题的关键：令含参项系数 $=0$，或化简后参数自动消去};

% 竖分割线
\draw[gray!50, dashed, thin] (2.7, 0.8) -- (2.7, 6.2);
\draw[gray!50, dashed, thin] (6.3, 0.8) -- (6.3, 6.2);

\end{tikzpicture}
\end{document}
